% !TeX root = ../papers/nc-p3.tex
\section{全有限窓の半安定性と厳密半安定枝の全次数正則性}
\label{sec:finite-window-semistability}
\begin{theorem}\label{thm:all-finite-windows-semistable}
四辺の点対応が有限な全ての標準有限Frobenius窓は半安定である。
厳密半安定、すなわち安定でない場合は全て標準quadratic Koszul左右AS正則で、
唯一の多重安定極限も標準正則である。特に有限点グラフにおいて全次数判定がなお必要な候補は、
閉軌道と有限安定群を持つ安定点に限られる。
\end{theorem}

\begin{proof}[Hilbert--Mumfordの全不等式]
四位置の任意の一径数群を対角化し、各位置の文字重みから最低値を引いて
$0,d_i,e_i,f_i$ の昇順とする。窓で比較すべき平均は $M=\sum_i(d_i+e_i+f_i)/4$ である。
辺 $i$ の六関係に低い重みから初期部を取り、その独立な初期基底の重みを
$r_{i,1}\le\cdots\le r_{i,6}$ とする。
最低文字の一点と隣の低い二文字の直線との積に関係が一つもなければ有限性に反する。
最高文字だけを片側に持つ関係は純テンソルになり存在しない。従って
\[
 r_{i,1}\le\min(d_i,d_{i+1}),\qquad
 r_{i,6}\le\min(e_i+f_{i+1},f_i+e_{i+1}).
\]
窓テンソルの最低重み $m$ は各 $k$ で $r_{i,k}+r_{i+2,7-k}$ 以下である。
実際、完全対合の係数行列で最初の $k$ 行と最初の $7-k$ 列の交叉が恒等零になることはない。
適合基底の初期テンソル積は各重み層で独立である。
$k=1$ の四本を平均し、二つの最小値をそれぞれ算術平均で抑えれば $m\le M$ を得る。
これが半安定性である。

等号なら用いた各不等式が等号である。特に $d_i=d$、$f_i-e_i=g$ は共通であり、
$r_{i,1}=d$, $r_{i,6}=e_i+e_{i+1}+g$ となる。
$m=d+e_i+e_{i+1}+g$ が全 $i$ で同じなので $e_i=e_{i+2}$ である。
したがって矢重みは $0,d,d+s,d+s+g$ と $0,d,d+t,d+t+g$ が交互に並ぶ。
$0\le s\le t$ としてよい。以下スカラーでない等号方向を調べる。
\end{proof}

\begin{proof}[すべての切断階数を保つこと]
まず $d=0$ は不可能である。$g>0$ なら最低空間の次元 $l_i$ は
中央差が正なら2、零なら3で、最高の次の空間の次元 $h_i$ はそれぞれ1,3である。
最低関係の本数は低い二固有空間の積への有限切断から $l_i+l_{i+1}-2$ 以上である。
これらは反対辺の最高関係としか対合できない。最高関係は最高文字と次高文字の二通りの積と
最高どうしの積だけに入り、両側の次高部分への射影が純関係の不在から単射である。
よって本数は $\min(h_i,h_{i+1})$ 以下である。$l$ の組 $2,2$、$2,3$、$3,3$ の全てで矛盾する。
$d=g=0$ なら非スカラー位置の重みは二重ずつ、スカラー位置の重複度は4である。
非自明な方向では各辺の最低空間の次元は $(2,2)$ 又は $(2,4)$ で、最高も同じ組である。
有限な最低切断に必要な本数は2又は4だが、最高ブロックでSegreを避ける六関係部分の
次元上界は1又は3である。これも両立しない。

これで $d>0$ であり、最低点どうしが点列をなす。三次元接点がある場合は最後に扱う。
それ以外ではこの点対で六関係の一次部は少なくとも四本独立である。
二端の重みを上の $s,t$ の順に並べると $M=3d+s+t+g$ である。
各辺で添字 $i$ を略し、次の上界が得られる。
\[
\begin{array}{ll}
 r_2\le d+\min(d,s), & r_5\le 2d+t+g,\\
 r_3\le d+s+\min(d,g), & r_4\le\min(d+t+g,\,2d+s+t).
\end{array}
\]
実際、低い $a,b$ 文字の射影空間の積への有限切断には $a+b-2$ 本以上必要である。
$r_2$ には $2,2$ と $3,1$（小さい中央差の側で3）、
$r_3$ には $4,1$ と $3,2$、$r_4$ には $3,3$ を使う。
$r_4$ のもう一つの上界は、$d+t+g$ を超える関係が最低点で一次部を持たないことによる。
$r_5$ の上界を超える項は最高二文字ずつの積だけに入る。ここで純テンソルを避ける空間は高々一次元である。

反対辺も同じ上界を持ち、$r_{i,2}+r_{i+2,5}\ge M$ から $s\le d$、
$r_{i,3}+r_{i+2,4}\ge M$ から $d\le g\le d+s$ を得る。
後者では $\min(d,g)+\min(g,d+s)\ge d+g$ を使った。
これらの範囲では各対の上界の和がちょうど $M$ である。
よって六重みは全辺で正確に
\[
 d,\ d+s,\ 2d+s,\ d+t+g,\ 2d+t+g,\ 2d+s+t+g
\]
となる。相補切断の行列式のどの非零項も六係数の重みの和が $6M$ で、
各非零係数の重みは $M$ 以上である。従って重み $M$ の部分だけでも行列式は非零であり、
極限窓の隣接二切断は階数6である。
\end{proof}

\begin{proof}[中央差の一致]
$s<t$ で $s<g$ とすると、重み $r_5$ 以上の全項のうち、$t$ 側の最高文字を含まないものは
もう一方の最高文字と $t$ 側の次高文字の積一つだけである。
実際次高どうしの重みは $2d+s+t<r_5$ で、最高と重み $d$ の組も $t$ 側が最高でなければ $r_5$ 未満である。
上位二関係をこの一成分へ射影した核には純関係が生じ、矛盾する。
従って $s<t$ なら $s=d=g$ である。単位を $d$ として矢重みは
$0,1,2,3$ と $0,1,1+u,2+u$ ($u>1$)、関係は $1,2,3,2+u,3+u,4+u$ で全て単純である。
文字を昇順に $0,1,2,3$ と呼ぶ。最低関係の $01,10$ は最低点の両評価からともに非零である。
最高関係の $23,32$ の初期係数も、どちらかが欠ければ最高どうしの補正だけが残って純関係となるため、ともに非零である。
極限の第一切断の階数1部分で $0132$ は非零で、次の切断には二文字 $13$（第二位置の重み1、第三位置の重み3）があり、
その和4が関係重みでなければならない。従って $u=2$ である。
このとき $30$（重み3）は最初の三関係のうち第三の初期部に必要である。
第二位置の最低点で評価した階数が3であり、第四以後の関係は第一位置の文字をそこへ供給できないからである。
$03$（重み4）は第四関係の初期部に必要である。第五以後は一次部が零であり、最低点のJacobi階数は4以上だからである。
ゆえに極限の $3003$ の係数がこの二つの単純重み間の完全対合により非零となる。
ところが次の切断には $00$ の関係がない。これで $s=t$ が従う。
\end{proof}

\begin{proof}[等号方向のPBW還元]
$s=t=0$ なら $g=d$ で、矢重みは $0,1,1,2$ にできる。重み1の両混合射影は、
各最低点での評価により二次元の同型である。重み3の両混合射影は
実際の最高二関係における純関係の排除から同型である。重み2では $03,30$ への射影が
最低点のJacobi階数4以上から同型である。命題\ref{prop:double-middle-finite-pbw}の証明で
極限窓のPBW性を得る。その証明で用いる中央の二行列の可逆性は、
重み1対3の完全対合の両混合射影が同型であることから従い、極限の点グラフの有限性は要らない。

$s=t>0$ の場合、四文字は全て単純である。重み $d$ の $01,10$ と最高重みの $23,32$ は
前段と同様に全辺で非零である。この二つの関係線が対辺で完全に対合することを別の切断に読むと、
$02,20$ と $13,31$ の四係数も各々の関係線で非零となる。
実際 $0132,1023$ の中央対と外対を使い、開始位置を全て回せばよい。
六つの関係重みのうち中央二つを除く四つは単純で、各関係の左辞書順先頭は $10,20,31,32$ である。
三つの不等式 $0<s\le d\le g\le d+s$ からこれらの重みで
あり得る余分な項は第二の $11$、第五の $22$ だけで、いずれも先頭を変えない。

$g>d$ なら中央二重みも単純で、$12,21$ と $03,30$ は $0123,1032$ を別切断で読むとともに非零である。
前者の先頭は $21$、後者では $22$ の追加が可能でも先頭は $30$ である。
$g=d$ なら中央二関係の $03,30$ への射影は最低点でのJacobi階数から同型である。
両対の $03,30$ 主小行列は $00,33$ の不在から可逆対角である。
行 $21$、列 $03$ の係数は別切断の $10,32$ により非零なので、
$30$ を除いた関係の先頭は $21$ となる。これで全場合で六つの下降対が先頭である。
各一文字の全支持はこれらの非零係数と単純重みから従う。
四テンソルの独立な四縮約が三次の重なりを全て解消し、定理\ref{thm:simple-weight-finite-windows}の
PBW証明により極限は標準Koszul左右AS正則である。命題\ref{prop:unrestricted-good-reduction}で元の窓にも戻る。

三次元接点があれば定理\ref{thm:maximal-tangent-ore}の純三文字関係が
ある点縮約の中央支持をも三次元にする。実際その縮約の末尾縮約像は
点での左評価核の三空間であり、これは接次元3により右評価核でもある。
定理\ref{thm:finite-cubic-separators}により元の窓は正則で、命題\ref{prop:maximal-tangent-closed-limits}で
多重安定極限も正則である。以上、三次元接点のない場合には任意の非自明な等号方向で上の証明が働く。
厳密半安定点では数値判定と閉軌道への一径数群判定により、必要な等号方向を選べる。
すでに多重安定で安定でない場合は正次元の簡約安定群内の一径数群を使う。
\end{proof}
